\documentclass[fleqn]{article}
\usepackage{amsmath}
\usepackage{amssymb}
\usepackage{array}
\usepackage[left=2cm,right=2cm,top=1.5cm,bottom=1.5cm]{geometry}
\setlength{\mathindent}{1.5em}

\newcommand{\Rp}{R\bigl(B_i^{\,2j+1},\, B_i^{\,2j},\, c_i^{\,j}\bigr)}
\newcommand{\Rt}{R_2\bigl(c_i^{\,2},\, B_i^{\,3}\bigr)}

\begin{document}

\begin{center}
{\LARGE \textbf{NR4SD$^{-}$}}\\[4pt]
{\large DP8 (8$\times$4) --- non-redundant signed digits of $B$ with an early carry, three per nibble}
\end{center}

\vspace{1.5em}

\noindent\textbf{Step 1 --- the dot product}
\[
\sum_{i=0}^{7} A_i \, B_i
\]

\vspace{0.5em}

\noindent\textbf{Step 2 --- recode $B$ into radix-4 digits and exchange the sums}

\noindent Each $B_i$ is a 4-bit value read as signed. Each pair plus its carry in gives a
partial sum; a sum of $2$ or more is pushed up as a carry, one step earlier than
NR4SD$^{+}$, which leaves a digit in $\{-2, -1, 0, 1\}$. As in NR4SD$^{+}$ the signed
read of the top bit and the last carry form one extra digit:
\[
s_i^{\,j} = 2\,B_i^{\,2j+1} + B_i^{\,2j} + c_i^{\,j}
\qquad
c_i^{\,j+1} = \bigl[\, s_i^{\,j} \geq 2 \,\bigr]
\qquad
d_i^{\,j} = s_i^{\,j} - 4\,c_i^{\,j+1}
\qquad
c_i^{\,0} = 0
\qquad
d_i^{\,2} = c_i^{\,2} - B_i^{\,3} \in \{-1, 0, 1\}
\]
\[
\sum_{i=0}^{7} A_i \, B_i
  \;=\; \sum_{i=0}^{7} \Biggl[\, \sum_{j=0}^{2} 4^{j} d_i^{\,j} \Biggr] A_i
  \;=\; \sum_{j=0}^{2} 4^{j} \sum_{i=0}^{7} d_i^{\,j} \, A_i
\]

\noindent Two digits span $[-10, 5]$ and miss $6, 7$, so a nibble needs the third digit;
when $B_i$ is a lower slice of a wider $B$ the top bit is read unsigned and
$d_i^{\,2} = c_i^{\,2}$.

\vspace{0.5em}

\noindent\textbf{Step 3 --- the partial products}
\[
\sum_{i=0}^{7} A_i \, B_i
  \;=\; \sum_{j=0}^{1} 4^{j} \sum_{i=0}^{7}
        \underbrace{\Bigl[\, d_i^{\,j} \, A_i \Bigr]}_{\textstyle \Rp}
  \;+\; 16 \sum_{i=0}^{7}
        \underbrace{\Bigl[\, d_i^{\,2} \, A_i \Bigr]}_{\textstyle \Rt}
\]

\noindent $\Rp$ depends on the pair and its carry in and takes one of four values;
$\Rt$ is the top digit:

\vspace{0.5em}

\begin{center}
\renewcommand{\arraystretch}{1.2}
\begin{tabular}{c c c | c | c | l}
$B_i^{\,2j+1}$ & $B_i^{\,2j}$ & $c_i^{\,j}$ & $c_i^{\,j+1}$ & $d_i^{\,j}$ & \multicolumn{1}{c}{$R$} \\
\hline
$0$ & $0$ & $0$ & $0$ & $0$  & $0$ \\
$0$ & $0$ & $1$ & $0$ & $+1$ & $A_i$ \\
$0$ & $1$ & $0$ & $0$ & $+1$ & $A_i$ \\
$0$ & $1$ & $1$ & $1$ & $-2$ & $-2A_i$ \\
$1$ & $0$ & $0$ & $1$ & $-2$ & $-2A_i$ \\
$1$ & $0$ & $1$ & $1$ & $-1$ & $-A_i$ \\
$1$ & $1$ & $0$ & $1$ & $-1$ & $-A_i$ \\
$1$ & $1$ & $1$ & $1$ & $0$  & $0$ \\
\end{tabular}
\qquad\qquad
\begin{tabular}{c c | c | l}
$c_i^{\,2}$ & $B_i^{\,3}$ & $d_i^{\,2}$ & \multicolumn{1}{c}{$R_2$} \\
\hline
$0$ & $0$ & $0$  & $0$ \\
$0$ & $1$ & $-1$ & $-A_i$ \\
$1$ & $0$ & $+1$ & $A_i$ \\
$1$ & $1$ & $0$  & $0$ \\
\end{tabular}
\end{center}

\vspace{0.8em}

\noindent\textbf{Step 4 --- split the work}

\noindent The encoding is a function of $B$ alone, so it is done in the dispatcher, which
sends the digits $d_i^{\,j}$ --- the control signals of the selection multiplexers --- to the
PE; only the selection of the multiple stays inside:
\[
\underbrace{\bigl\{\, B_i^{\,2j+1},\, B_i^{\,2j},\, c_i^{\,j} \,\bigr\} \;\longrightarrow\; d_i^{\,j}
  \qquad
  \bigl\{\, c_i^{\,2},\, B_i^{\,3} \,\bigr\} \;\longrightarrow\; d_i^{\,2}}_{\textstyle \text{dispatcher: encoding}}
\qquad\longrightarrow\qquad
\underbrace{\sum_{j=0}^{2} 4^{j} \sum_{i=0}^{7} d_i^{\,j} \, A_i}_{\textstyle \text{PE: selection}}
\]

\noindent Each $d_i^{\,j}$ travels as a 2-bit code which the cell uses directly: 6 control
bits per lane, a carry chain in the encoder, none in the cell.

\vspace{0.8em}

\noindent $B$ is the narrow operand: 3 digits $\times$ 8 lanes = 24 selections of
10 bits per DP8, each over $\{0, A_i, -A_i, -2A_i\}$ --- the shifted multiple only appears
negated, so the shift and the negation always combine; carries fire on $2$ and $3$, more
often than in NR4SD$^{+}$.

\end{document}
